\part*{Lecture 5: Black's Formula and Callable Bonds}
\addcontentsline{toc}{part}{Lecture 5: Black's Formula and Callable Bonds}
\beginlecture{05}

Lecture~5 is the first session of FINM~37400 in name, but the lecturer
treated it as a continuation of FINMAT~374: ``for those of you in both
courses, this is really going to feel like you did the pre-midterm work,
you did the midterm, now we're doing the post-midterm work up to the
final.'' The single 147-minute session does roughly four jobs: it
\emph{closes out} the swap discussion that ran out of clock in
Lecture~4 (notional vs.\ gross vs.\ net market value, the real drivers
of swap demand, and a short LTCM-flavoured revisit of the swap-spread
trade); it \emph{transitions} to options on rates and bonds without
re-doing FINMAT~330's Black--Scholes derivation; it \emph{works through}
Black's formula and the pricing of a callable Freddie~Mac bond,
introducing the option-adjusted spread (OAS) along the way; and it
\emph{previews} caps, floors, and swaptions, which the lecturer flagged
as the actually-important products of fixed-income derivatives. The
notebooks in scope are \emph{6.1.~Black's Formulas}, \emph{6.2.~Callable
Bonds}, and the appendix \emph{6.X.1.~Note on Black-Scholes} (referenced
but not lectured on --- ``go back to your FINMAT~330 notes if you want
to derive it''). The cap-and-floor primer at the end is taught from
slides that belong to the next chapter (7.1.~Forward Volatility and
7.2.~Swaptions); these notes preserve the order in which the lecturer
presented them.

\section{Wrapping up swaps from FINMAT~374}\label{L5-sec:swaps-wrap}

The lecturer opened by noting that swaps ``straddle both courses'' ---
they were started in FINMAT~374 but the punchline was deferred. Two
loose ends needed tying off: \emph{why} the swap notional is so
enormously larger than the actual capital at risk, and \emph{why} the
swap market is the most liquid rate market we will see, larger even
than the Treasury market. The first is a definitional question; the
second is a story about hedging duration.

\subsection{Notional, gross market value, and net}\label{L5-sec:swaps-notional}

\begin{keyblock}
\textbf{Slide L5.1.} \emph{(How big is the swap market?)}\\
Total notional outstanding on interest-rate swaps is on the order of
\$400~trillion. This is \emph{not} \$400~trillion of wealth tied up
in the asset class.
\begin{itemize}
  \item \emph{Notional} is the face amount used to scale the cash flows.
        It is never exchanged.
  \item \emph{Gross market value} is the sum of the absolute mark-to-market
        values across all positions. For interest-rate swaps, this number
        is roughly \$10~trillion --- a few percent of notional.
  \item \emph{Net market value} is gross net of offsetting positions held
        by the same counterparty.
\end{itemize}
\end{keyblock}

The lecturer warmed up with a thought experiment: ``if you told a
family member there's \$400~trillion of wealth in swaps, they would
say, you fell asleep in class.'' The reason that headline figure is so
misleading is that the dollar exposure on a single swap is
\emph{zero} the moment the contract is signed --- the fixed leg is set
exactly so that the swap has no entry cost. Only after rates move does
either leg accrue value. On a \$100~million 10-year swap, even a
substantial rate move builds maybe a percent or two of dollar value per
year of remaining life --- ``maybe \$10~million on \$100~million
notional, certainly nothing close to the notional itself.'' That is
how \$400~trillion notional becomes order \$10~trillion gross.

The further compression from gross to net comes from how swap
positions are managed in practice. The lecturer used an in-class
example: a student who is up \$10~million on a swap with the lecturer,
but down \$7~million on an offsetting swap with a third student, holds
the same notional of \$100~million but a \emph{net} value of around
\$3~million. ``These contracts are not typically created and
destroyed; when you want to get out, you create an offsetting swap.''
The headline notional inflates as the population of contracts grows,
but the net dollar value grows much more slowly.

\begin{checkpoint}
\textbf{My note.} Even after every netting adjustment, the swap market
dwarfs the equity-options market on the same gross-vs-net measure ---
``\$400~trillion notional, \$10~trillion gross, some smaller fraction
net, but still much, much bigger than the netted value of US equity
call options.'' That asymmetry is what makes swaps the dominant
rate-derivatives market. The reasons \emph{why} are the subject of the
next subsection.
\end{checkpoint}

\subsection{Why is the swap market so large? Duration management}\label{L5-sec:swaps-duration}

The naive story --- corporates issue floating but want fixed, find an
opposite-handed corporate, and the two swap --- is real but not the
main driver. ``It seems hard to believe that's driving what is
essentially the biggest market we've seen yet.'' The lecturer pressed
the room for the missing story. One student offered \emph{speculation}
on rates, which the lecturer accepted but redirected; the canonical
answer is closer to the rates-trader analog of equity-beta hedging.

\begin{keyblock}
\textbf{Slide L5.2.} \emph{(Why swaps are huge.)} The dominant driver
of swap demand is \emph{level-of-rates risk management} --- the
fixed-income analogue of equity-beta hedging.

\textbf{Equity analogy.} An equities portfolio manager runs many
long/short positions she has views on. She insists her bets are
about specific stock pairs, not the market. The risk manager replies:
``you've still taken on a lot of beta. Keep doing your thing; I will
quietly overlay an S\&P-future hedge to neutralize broad-market
exposure.''

\textbf{Rates analog.} A rates trading desk runs many bond positions
based on relative-value views. Incidentally, those positions inherit
\emph{duration} --- exposure to the level of rates. The risk manager
will use \emph{swaps} to neutralize that duration without disturbing
the underlying trades.
\end{keyblock}

The lecturer then unpacked why a swap is uniquely good at this. A
\emph{receiving-fixed} 10-year swap mimics owning a 10-year Treasury:
the duration of the fixed leg is roughly the modified duration of a
fixed-coupon bond with the same maturity (around 7--8 for a 10-year).
The floating leg on the other side has near-zero duration --- the
floating-rate-note result from Lecture~2 says that the cash flows
adjust to rate moves and are simultaneously discounted at moving rates,
and the two effects cancel almost exactly. So a fixed-for-floating
swap is essentially a long-duration leg paired with a zero-duration
leg, and trading the swap moves duration the same way trading the
underlying Treasury would.

\begin{tipblock}
\textbf{Filling the gap.} The lecturer's claim ``the floating-rate
note has no duration'' is the punchline of \emph{2.4~Floating-Rate
Notes}. A floating-rate note pays $r_t \cdot \tau$ each period and
returns face at maturity. Each coupon is set in arrears by the same
short rate that discounts it back, so a parallel shift in the term
structure leaves both numerator and denominator scaled identically and
the present value unchanged. This is exactly the duration-equals-zero
result; it does not require taking derivatives explicitly.
\end{tipblock}

The lecturer's punchline was that swaps then do duration-management
\emph{on better terms} than the equivalent Treasury short would. A
short Treasury position must be repo'd for years and consumes balance
sheet through the haircut and the rolled financing. A swap, in
contrast, is no-cost to enter: ``I'm offloading a massive amount of
duration by promising her \$100~million swap, but how much did it take
on my balance sheet? Nearly nothing, because I'm not paying her
\$100~million.'' Margin posted to the central counterparty is small,
much like a futures margin account. That balance-sheet efficiency is
what makes swaps the rate-hedger's instrument of choice and the reason
the swap market dwarfs the cash-Treasury market on the
DV01-scaled-volume measure introduced in Lecture~4.

\begin{checkpoint}
\textbf{My note.} The lecturer ranked swaps ``closer to the order of
Treasuries than on the order of looking at bond options.'' For
intuition: when sizing the toolkit of fixed-income products, swaps
belong with the underlying-rate instruments (Treasuries, repo, STIR
futures), not with the derivatives layer above them. They are big
enough and deep enough that \emph{everything else} is priced relative
to a swap curve.
\end{checkpoint}

\subsection{Revisiting the swap-spread trade}\label{L5-sec:swap-spread}

Before moving to today's content, the lecturer reopened the
\emph{5.2~Swap-Spread Trade} case from the FINMAT~374 homework, both
because part of the class was new to it and because the
fixed-income-derivatives course will return to it through swaptions,
forward-vol trades, and CTD-like relative-value structures.

\begin{keyblock}
\textbf{Slide L5.3.} \emph{(Swap-spread trade in 2008.)} The
30-year US Treasury at par pays a 4.5\% coupon. The 30-year swap rate
is 4.26\%. The cash-flow trade:
\[
  \text{long Treasury (receive 4.50\%) } + \text{ pay-fixed swap (pay 4.26\%)}.
\]
Repo finances the long Treasury at $r_{\text{repo}}$, and the swap
floating leg pays $r_{\text{repo}}$ approximately. The two financing
flows cancel, leaving a net carry of $4.50\% - 4.26\% = 24$~bps per year.
At 30$\times$ leverage on a 2\% repo haircut, that becomes roughly
$24~\text{bps} \times 30 \approx 7.5\%$ on equity.
\end{keyblock}

The lecturer's first question to the class was the obvious one:
\emph{is this an arbitrage}? In the strict sense, no --- the position
is not zero-state at every horizon. But over the full 30-year life,
abstracting from credit risk on the clearing party and the
convexity-adjustment frictions of the daily-settled swap leg, the two
``risk-free'' rates at maturity must coincide; the spread is
mechanically forced to converge. ``By assuming away these frictions, I
have two different risk-free providers for cash, and cash has to have
the same price.''

The much more interesting question is over what \emph{horizon} the
trade pays off. The lecturer pushed: ``is it overnight? a week? a
month? Absolutely not --- there's no mathematical law forcing the
spread to converge tomorrow.'' And once you put on 50$\times$ leverage,
intermediate margin calls become the dominant risk. He walked the
class through the LTCM-style death spiral: at 24~bps the trade looks
boring, so you lever it; once levered, an adverse move in the spread
generates margin calls big enough to consume the equity capital that
sized the trade in the first place. ``If your long-treasury position
is down, your repo counterparty will say the haircut isn't big enough
anymore --- you have to put up more.''

\begin{tipblock}
\textbf{Filling the gap.} The lecturer asked the class \emph{why} the
swap spread has stayed negative for 15 years, contrary to the
arbitrage intuition. His answer pulled from Boyarchenko (Federal
Reserve) and the post-2009 regulatory regime. Bank prop-trading desks
had the balance sheet to attack such trades and were the natural
arbitrageurs; Dodd-Frank/Volcker reduced their appetite. The
remaining capital --- hedge funds, prop traders --- prices in the
financing risk and the indeterminate horizon, and the trade is not
attractive at the size that would tighten the spread. This is the
``limits to arbitrage'' phenomenon Vrugt-style spread trades exhibit
generally: the math says converge, the institutions say it depends on
funding.
\end{tipblock}

\begin{checkpoint}
\textbf{My note.} The swap-spread case is also a sneak peek at why
\emph{forward volatility} (Lecture~6) matters: a leveraged
relative-value trader cares about the path between today and
convergence, not just the terminal payoff. Forward-vol surfaces tell
you what option-implied movement to expect at \emph{intermediate}
horizons. The swap-spread trade is essentially a directional bet that
those intermediate movements will not consume the equity --- a bet
that some forward distribution is tighter than the market thinks.
\end{checkpoint}

\section{From spot options to options on rates}\label{L5-sec:transition}

\begin{keyblock}
\textbf{Slide L5.4.} \emph{(Today: options on rates and bonds.)} The
remainder of the course --- and the entirety of this lecture from this
point on --- addresses derivatives whose underlying is a rate, a
forward rate, or a bond price. The plan is \emph{not} to redo
FINMAT~330's Black--Scholes derivation; that course remains the
prerequisite. Today's job is to identify what breaks when the
underlying is fixed-income, fix it with the smallest possible
modification (Black's model), and price one example end-to-end (the
Freddie~Mac callable bond).
\end{keyblock}

The lecturer's framing was deliberate. FINMAT~330 covered the
classical Black--Scholes formula in depth: replication, the partial
differential equation, the greeks, and the risk-neutral expectation. A
student who took it should be able to use the formula, recognize the
implied-volatility input, and reason about a Greek. The fixed-income
derivatives course assumes that comfort and treats the formula as a
\emph{starting point} from which two things must be patched: (1) the
risk-free rate is not constant, and (2) the underlying does not look
remotely like a Brownian motion when the underlying is a bond.

The lecturer also signalled what is being deferred. The \emph{6.X.1
Note on Black-Scholes} appendix on the live page contains a
self-contained derivation --- Brownian motion, geometric Brownian
motion for the stock, Itô's lemma, the replicating portfolio, the
fundamental PDE, the greeks, the theta--gamma identity, and the
risk-neutral expectation form. He was explicit: ``go back to your
FINMAT~330 notes if you want to derive it; I am not going to spend
two more hours on it to prove the forward measure exists.'' The
appendix is on the live notebook for self-study; it is not part of
what the lecture covered.

\subsection{Why callable bonds, when they are the least important
product?}\label{L5-sec:why-callable}

The lecturer was up-front that callable bonds will not be a major
working object of the course. ``Callable bonds are probably the least
important derivative we are going to talk about.'' The motivation for
spending the day on them was pedagogical: they are the \emph{minimum
modification} of a FINMAT~330 problem that exposes every fixed-income
wrinkle worth seeing.

\begin{checkpoint}
\textbf{My note.} The hierarchy of importance for the rest of the
course, in the lecturer's ranking:
\begin{enumerate}[label=(\arabic*)]
  \item \emph{Swaps} (already covered in FINMAT~374; the dominant rate
        derivative).
  \item \emph{Treasury futures} (covered in Lecture~4 alongside the
        cash market).
  \item \emph{Caps and floors} --- ``a monster in terms of size
        compared to many of these other things you've heard of.''
  \item \emph{Swaptions} --- ``one of our big three.''
  \item \emph{Callable bonds} --- the gateway, ``a perfect
        intersection of those worlds.''
\end{enumerate}
The order matters: caps/floors and swaptions are the marquee
fixed-income derivatives, but pricing them requires forward-volatility
machinery that does not appear until Lecture~6. The callable bond is
the simplest object that requires Black's formula plus a
fixed-income-aware vol input, so it is the right exercise for today.
\end{checkpoint}

\section{Callable bonds: the embedded short call}\label{L5-sec:callable-setup}

\begin{keyblock}
\textbf{Cells~6.2.0--9.} \emph{(Page 6.2 \textperiodcentered{} Callable
Bonds.)}\\
\textbf{Definition.} A callable bond is a vanilla bond plus an
embedded call option held by the \emph{issuer}: at certain dates the
issuer may pay the strike (typically face value) and reclaim the bond.
The bondholder is therefore \emph{short} the call. Equivalently
\[
  V^{\text{callable}}_t \;=\; V^{\text{vanilla}}_t \;-\; C_t,
\]
where $V^{\text{vanilla}}_t$ is the price of an otherwise identical
non-callable bond and $C_t$ is the value of the embedded call option
on the bond, struck at face.
\end{keyblock}

The student answer the lecturer harvested in the room captured the
intuition: ``a callable bond is one the issuer has an option to recall
at face value; usually the issuer calls when they can refinance later
at a lower rate.'' The same student noted negative convexity --- the
issuer calls precisely when the bond would otherwise be valuable, so
the holder's upside is capped while the downside is unchanged. The
lecturer translated both into FINMAT~330 language: \emph{negative
convexity} on the bond is the same statement as \emph{negative gamma}
on the embedded call once the bond's own positive convexity is
overwhelmed.

\begin{tipblock}
\textbf{Filling the gap.} The decomposition $V^{\text{callable}} =
V^{\text{vanilla}} - C$ is the entire pricing strategy and deserves
unpacking. The issuer's call, evaluated at the bond's clean price,
pays the issuer $\max(P_{\text{bond}} - K, 0)$ on exercise. The
bondholder receives the lesser of $P_{\text{bond}}$ and $K$ at the
call date --- their position is $P_{\text{bond}} - \max(P_{\text{bond}}
- K, 0) = \min(P_{\text{bond}}, K)$. Subtracting a long call from a
long bond is the same algebra. To value the callable, value the
vanilla and value the call separately; the work shifts entirely to
producing the call value via Black.
\end{tipblock}

\subsection{Bermudan exercise}\label{L5-sec:callable-bermudan}

The room's first guess was that the callable's exercise window is
``European'' or ``American.'' Both are wrong. Freddie~Mac callables ---
and most agency callables --- are \emph{Bermudan}: exercisable on a
discrete set of pre-specified dates, typically every six months
beginning some lockout period after issue. ``Every time the agency
pays a coupon, they have the choice of whether to call the bond back
in.''

The lecturer emphasised \emph{when} the Bermudan-ness matters and when
it doesn't. There are two regimes:
\begin{enumerate}[label=(\arabic*)]
  \item \emph{Long lockout, short remaining life} --- e.g.\ a 10-year
        bond first callable in 7 years. From issue to first-call this
        looks essentially European; the embedded option is exercised
        only if the lockout-end value is above strike.
  \item \emph{Short lockout, long remaining life} --- e.g.\ a 15-year
        bond first callable in 6 months. There are 30+ semi-annual
        exercise dates; this is much closer to American. ``If I had to
        approximate one or the other, I would approximate it as
        American.''
\end{enumerate}
The Freddie~Mac bonds the case study uses are in regime~(2). The
lecturer's plan for today was to ignore the Bermudan structure
entirely and price the call as European with strike $K=100$ and
expiration at the first call date; the case study and exercise revisit
that approximation numerically. The justification, beyond
convenience, is that exercising a Bermudan call \emph{early} discards
the remaining optionality, so the issuer typically waits unless the
bond is deeply in the money or close to maturity.

\begin{checkpoint}
\textbf{My note.} Why \emph{not} exercise an American/Bermudan call
early? Same reason a stockholder does not exercise an American call on
a non-dividend-paying stock: cashing out terminates the gamma. The
issuer will exercise only when the remaining gamma is small relative
to the dollar gain from calling. Translating the FINMAT~330
heuristic: vol low, time-to-maturity short, deeply in the money. For a
bond with 15 years remaining and the dirty price still near par, none
of those conditions hold, and the European approximation should be
tight.
\end{checkpoint}

\subsection{Why a Treasury proxy: Freddie Mac}\label{L5-sec:freddie}

Most callable corporate bonds carry the issuer's credit risk on top of
the embedded call, and credit modeling is explicitly out of scope for
this course (it lives in Alex Popovici's spring class). The
\emph{agency} callables --- Freddie~Mac, Fannie~Mae, Ginnie~Mae --- are
the workaround. The agencies were absorbed into the Treasury after the
2008 financial crisis and remain backed by the full faith and credit
of the US government. ``Even though it's not issued as a Treasury
bond, it should seemingly have identical credit risk because it
ultimately is the same issuer.''

\begin{keyblock}
\textbf{Cells~6.2.2--6.} \emph{(Why Freddie Mac.)}\\
\textbf{Freddie Mac (Federal Home Loan Mortgage Co.).} A US housing
agency, taken into the Treasury after 2008. For pricing purposes,
treat its debt as credit-risk-free (Treasury-equivalent issuer).

\textbf{Spread to Treasuries.} Freddie~Mac bonds typically trade about
20~bps higher in yield than maturity-matched Treasuries. This is real
but small; it reflects mostly liquidity (small issue size, irregular
issuance, less global attention) rather than credit.

\textbf{Why use them in this course?} The Treasury itself does not
issue callable bonds today. To get a clean callable instrument
\emph{without} corporate credit risk, agency callables are the only
choice.
\end{keyblock}

\subsection{The bond table the lecturer used}\label{L5-sec:freddie-table}

\begin{keyblock}
\textbf{Cell~6.2.7.} \emph{(The Freddie Mac bonds in scope.)}\\
Quotes from February~2024. Three bonds, all issued in
January--February~2024 (so essentially fresh issues at quote date).
\begin{itemize}
  \item Maturities in 2034 and 2039 --- a 10-year and a 15-year.
  \item First call date is six months after issue. Subsequent call
        dates every six months thereafter.
  \item Coupons are 5\%, 8\%, and 6\% --- fixed, paid semi-annually.
  \item Issue sizes are tiny relative to Treasury sizes; liquidity
        will not match the deep cash-Treasury market.
\end{itemize}
\end{keyblock}

The lecturer pointed at the dirty price column and noted that the most
recently-issued bond traded essentially at par on the quote date,
which makes the embedded call \emph{at-the-money} from the start ---
the strike is~100, the price is near~100, and there are 30
semi-annual exercise dates ahead. That arrangement should produce a
substantial call value. The case study notebook has the reader walk
through the actual valuations; today's job is to set up the machinery.

\section{Why Black--Scholes will not work directly}\label{L5-sec:bs-fails}

With the instrument defined, the lecturer turned to the modeling
question: can the call option on the bond be priced with the
classical Black--Scholes formula? The class diagnosed two failures.

\subsection{Constant rates is absurd in fixed income}\label{L5-sec:bs-fail-1}

The classical Black--Scholes model assumes a constant risk-free rate
$r$ that appears (a)~as the growth rate inside $d_1$ and (b)~as the
discount factor $e^{-r(T-t)}$ in front of the formula. In equity-options
land that assumption is a moderate inconvenience that practitioners
patch up with a flat term-structure approximation. In fixed-income
land it is fatal: the entire interest of the problem is that the rate
\emph{is} stochastic. ``If the interest rate is fixed and the same at
all maturities, we can just end class right now.''

\subsection{Pull to par: bonds are not Brownian}\label{L5-sec:bs-fail-2}

The deeper failure is that a bond price is \emph{predictable} in a
way that no stock price is. The class converged on the term ``pull to
par.'' At maturity the bond pays exactly~\$100 of face, so the bond
price as $t \to T_{\text{bond}}$ has a hard pin. Anywhere away from
maturity, the bond price will be \emph{drifting} toward~\$100 from
above (when rates have fallen and the bond trades over par) or from
below (when rates have risen and the bond trades under par).

\begin{keyblock}
\textbf{Slide L5.5.} \emph{(Pull to par, illustrated.)} The 30-year
Treasury issued in 1993 was originally priced at par. Rates fell
steadily through the 2000s, and by 2012 the bond traded around
\$160. By maturity in 2023 the price must return to \$100. So a
holder in 2012 knew with certainty that the price would fall by about
\$60 over the next 11 years --- ``is your best prediction of the future
price~\$160? It's impossible.''

The reverse case (Silicon Valley Bank, 2023) was just as predictable
in the other direction: SVB held Treasuries issued at 2\% coupons,
rates jumped to 4.5\%, the bond prices fell well below par. Their
expected path was \emph{up}; not random.
\end{keyblock}

\begin{tipblock}
\textbf{Filling the gap.} Why is this the killer for the original
Black--Scholes derivation? Black--Scholes pricing does not depend on
the stock's expected return $\mu$, by replication. Equivalently, in
the financial-economics view: a stock with high expected return has
high risk and is discounted at a correspondingly high rate; the two
exactly offset, and $\mu$ does not appear. For a bond, the expected
\emph{price path} comes from the difference between the bond's coupon
and the prevailing rate (the carry), \emph{not} from a risk premium
that is balanced by a discount adjustment. Concretely, a bond at
\$160 with an 8\% coupon when the risk-free rate is 5\% has expected
price-return roughly $-(8\%-5\%) \cdot \text{years remaining}/\text{remaining
amortization} \approx -5\%$ per year, while the holding period total
return is around the risk-free rate. The two no longer cancel ---
``the expected price path is not its expected return in a way that's
going to cancel out in Black--Scholes.''
\end{tipblock}

\begin{checkpoint}
\textbf{My note.} The lecturer's diagnosis is that Black--Scholes
needs an underlying whose expected drift is exactly absorbed by the
discount-rate machinery, leaving only volatility. Bond \emph{prices}
have a drift coming from carry that is not absorbed. Bond \emph{yields}
do better but introduce their own awkwardness; the lecturer is going
to use a third option --- bond \emph{forward prices} --- which absorb
carry by construction. The next section is that idea.
\end{checkpoint}

\section{Black's model: model the forward price}\label{L5-sec:black76}

The fix is the same one used in Black~(1976) for commodity futures: do
not model the spot price; model the \emph{forward} price. A forward
price already incorporates the carry between today and the forward
date, so it does not have the predictable drift that broke
Black--Scholes for bonds. Under the appropriate \emph{forward
measure}, the forward price is a martingale and the same Brownian
machinery used in FINMAT~330 carries through.

\subsection{From the spot formula to the forward formula}\label{L5-sec:black-classic}

\begin{keyblock}
\textbf{Cells~6.1.9--16.} \emph{(Black's Model -- Constant Rates.)}\\
Under \emph{constant} rates and lognormal forward prices, the European
call price on a forward-priced underlying is
\[
  c_t \;=\; e^{-r(T-t)}\,\Bigl[\, F_t\,\cN(d_1) \;-\; K\,\cN(d_2)\,\Bigr],
\]
\[
  p_t \;=\; e^{-r(T-t)}\,\Bigl[\, K\,\cN(-d_2) \;-\; F_t\,\cN(-d_1)\,\Bigr],
\]
where
\[
  d_1 \;=\; \frac{\ln(F_t/K) + \tfrac{1}{2}\sigma^2\tau}{\sigma\sqrt{\tau}},
  \qquad d_2 \;=\; d_1 - \sigma\sqrt{\tau},
  \qquad \tau = T - t.
\]
The forward price $F_t$ replaces the spot $S_t$ throughout. The drift
in $d_1$ is just $\tfrac{1}{2}\sigma^2$ --- not $r + \tfrac{1}{2}\sigma^2$
--- because the forward price already absorbs the cost-of-carry.
\end{keyblock}

\begin{tipblock}
\textbf{Filling the gap.} ``Where did the $r$ go in $d_1$?'' The
Black--Scholes $d_1$ has $r + \tfrac{1}{2}\sigma^2$ because the
risk-neutral drift of the spot price $S_t$ is $r$ (minus dividends).
The forward price has zero risk-neutral drift under its own measure
--- it is already the price of the future under no-arbitrage. So when
you re-express the formula in terms of $F_t$, the $r$ in the drift
disappears, and only the $\tfrac{1}{2}\sigma^2$ from the lognormal
correction remains. The discount factor $e^{-r(T-t)}$ in front does
not change.
\end{tipblock}

The lecturer pressed the class to spot the difference between this and
the textbook Black--Scholes. ``Can you see? In the old formula, what
was the top line? $S\cN - K^{*}\cN$ where $K^{*}$ is the present-valued
strike. Now it's $F\cN - K\cN$, and the entire bracket is multiplied by
$e^{-rT}$.'' The bracket-and-discount form is the practical
arrangement: pull the discount out so the bracket is the
\emph{undiscounted} expected payoff under the risk-neutral (here:
forward) measure.

\subsection{Why this is what we want for bonds}\label{L5-sec:black-bond-intuition}

The key sanity check is that the forward price of a bond is not pulled
to par the way the spot price is. The forward price satisfies the
carry identity from Lecture~4:
\[
  F_t \;=\; \frac{P_t - \sum_{T_i \in (t,T]} c_i Z(t,T_i)}{Z(t,T)},
\]
which already prices in the coupons collected and the financing accrued
between $t$ and $T$. Once those are stripped, the residual $F_t$ has no
deterministic drift and is much more nearly Brownian than $P_t$ is.
``On the forward price we are not going to have the issue that it's
just going down, down, down.''

For the Freddie~Mac problem the lecturer reminded the class \emph{which}
forward they need: not the forward to bond maturity (where there is no
uncertainty), but the forward to the option's expiration. ``The bond
matures in 15~years; the option expires whenever the option expires.
If the embedded option expires in 10~years, I need the forward I can
lock in today to buy a 15-year bond 10~years from now.'' That is the
$P_{\text{fwd}}(t, T_{\text{opt}}, T_{\text{bond}})$ object from
Lecture~4's forward-pricing formula --- with the option-expiration date
in the middle slot. Building it requires the discount factors $Z(t,
\cdot)$ from Lecture~1 and the bond's coupon schedule.

\subsection{Black's general formula: stochastic rates}\label{L5-sec:black-general}

The Black-1976 form above still has the constant-$r$ pathology in front
of the bracket. The next step --- and the technical core of the
lecture --- is to allow $r$ to be stochastic without losing the
clean expectation structure.

\begin{keyblock}
\textbf{Cells~6.1.27--35.} \emph{(Black's General Formula -- Stochastic
Rates.)}\\
\textbf{Switch numeraire.} Choose the zero-coupon bond $Z(t,T)$ as
numeraire. Under the corresponding measure (the \emph{$T$-forward
measure}, $\E^T$), the ratio $f(F_t,t)/Z(t,T)$ is a martingale for any
function $f$ of the forward price.

\textbf{Pricing identity.} The value of any derivative paying $H(F_T)$
at $T$ is
\[
  V_t \;=\; Z(t,T)\,\E^T_t\!\bigl[H(F_T)\bigr].
\]
For a European call $H(F_T) = (F_T - K)^+$ with $F_T$ lognormal under
$\E^T$, this reduces to
\[
  c_t \;=\; Z(t,T)\,\Bigl[\, F_t\,\cN(d_1) \;-\; K\,\cN(d_2)\,\Bigr],
\]
\[
  p_t \;=\; Z(t,T)\,\Bigl[\, K\,\cN(-d_2) \;-\; F_t\,\cN(-d_1)\,\Bigr],
\]
with the same $d_1, d_2$ as before:
\[
  d_1 \;=\; \frac{\ln(F_t/K) + \tfrac{1}{2}\sigma^2\tau}{\sigma\sqrt{\tau}},
  \qquad d_2 \;=\; d_1 - \sigma\sqrt{\tau}.
\]
The constant-$r$ discount $e^{-r(T-t)}$ has been replaced by the
observable, term-structure-dependent discount factor $Z(t,T)$.
\end{keyblock}

\begin{tipblock}
\textbf{Filling the gap.} The change of numeraire is the one piece of
this lecture the lecturer explicitly outsourced: ``go back to your
FINMAT~330 notes if you want to derive it.'' The intuition: the
risk-neutral expectation form $V_t = \E^Q[e^{-\int_t^T r_s\,\dd s}
H(F_T)]$ has the random discount factor \emph{inside} the expectation.
Because $r$ and $F$ are correlated, that integral cannot be pulled
out. Re-expressing the price as $V_t = Z(t,T) \E^T[H(F_T)]$ via
Girsanov's theorem and the change-of-numeraire formula extracts the
discount $Z(t,T)$ as a deterministic factor, leaving an expectation
over only the forward price. The extracted $Z(t,T)$ is the
\emph{price of a zero-coupon bond}, which the curve-fitting machinery
of Lecture~1 already produces. Operationally that is a strict win: the
hard input becomes a pre-existing course deliverable.
\end{tipblock}

The lecturer's verdict on this whole development was self-mocking:
``this is going to win the prize for the longest discussion you ever
had in financial math for the smallest amount of changes to the
equation. So I'm confident we just won that trophy.'' The Black~1976
classic and the Black general formula differ in exactly one place:
$e^{-r(T-t)}$ becomes $Z(t,T)$. Everything else is identical to the
formula from FINMAT~330. The justifying machinery (forward measure,
Radon--Nikodym, martingale property of $F_t/Z(t,T)$) lives in
FINMAT~330 and FINMAT~320; for this course it is enough to know what
to plug in.

\begin{checkpoint}
\textbf{My note.} The substitution $e^{-r(T-t)} \to Z(t,T)$ also tells
us \emph{which} discount curve we should use to price the call.
Whichever curve we use to bootstrap $Z(t,T)$ \emph{is} the curve we
implicitly trust as the risk-free rate in the model. For Treasuries
that's the on-the-run Treasury curve; for swap-related products
(coming up) the analogous choice is the swap curve. Different choices
of curve give different OAS values for the same option price; that
mismatch is one of the mechanisms behind a non-zero OAS, which we will
discuss shortly.
\end{checkpoint}

\section{Implied volatility, treated as a language}\label{L5-sec:ivol}

With Black's general formula in hand, the only remaining input is
$\sigma$. The lecturer used the implied-volatility discussion to make
two points: a methodological one (how implied vol is conceptually
analogous to YTM), and a practical one (what to plug in for our
Freddie~Mac problem when we do not have an option-pricing engine).

\subsection{Implied vol is to options what YTM is to bonds}\label{L5-sec:ivol-language}

\begin{keyblock}
\textbf{Slide L5.6.} \emph{(Implied vol as a quoting language.)}
The implied volatility of a market-traded option is
\[
  \sigma^{\text{imp}} \;=\;
   \bigl\{\, \sigma : C^{\text{Black}}(F_t, K, T-t, Z(t,T), \sigma)
              = C^{\text{market}} \,\bigr\}.
\]
That is, $\sigma^{\text{imp}}$ is the unique number that makes Black's
formula reproduce the market price. It is \emph{not} the realized
volatility of the underlying, and it is \emph{not} the model's
forecast of future volatility per se. It is a quote, like YTM is a
quote --- a single number that translates a dollar price into a
comparable scale across strikes and expirations.
\end{keyblock}

The lecturer drew the analogy explicitly. A yield to maturity is
neither a return nor a discount rate; it is a re-scaled price that
makes two bonds with different coupons and maturities directly
comparable. A 4\%-coupon bond at \$140 dirty and a 1\%-coupon bond at
\$98 dirty cannot be directly compared by dollars, but their YTMs put
them on a single axis. Implied volatility plays exactly the same role
for options. Two call options on the same underlying with different
strikes and different expirations cannot be directly compared by
premium, but their implied vols can. ``The 30-cent option and the
70-cent option --- which is the better deal? You'd have no idea. But
the 18\% implied-vol option vs the 35\% implied-vol option on the same
underlying tells you a story.''

The deepest point is that implied vol is computed using \emph{Black's
formula} regardless of whether the market actually uses Black to
price. American and Bermudan options are quoted in Black-implied vol
even though Black's derivation assumes European exercise. The quant
group inside a bank may use a Lévy-process model with stochastic
volatility and jumps; the trader will read its output, compare it
against the Bloomberg Black-implied-vol screen, and make a decision.
``Black is the universal translator.''

\subsection{What to plug in for our Freddie Mac call}\label{L5-sec:ivol-plug}

The intellectually honest workflow for valuing the Freddie~Mac
embedded call would be (a) build a serious bond-option pricing model,
(b) extract the model's implied vol on liquid quoted bond options, (c)
feed that vol into Black for the specific instrument. The lecturer
made clear we are not going to do (a) this week. ``We are going to
plug in a realized vol for the implied vol because we don't actually
have an engine yet to build a valuation of the option.''

That choice forces a unit conversion. The realized vol that is easiest
to estimate from data is the volatility of the bond's \emph{yield},
not the volatility of its \emph{forward price}. The case study quotes
``the realized vol on its YTM is roughly 2.68\% per year'' for the
underlying. Black's formula expects the vol of the forward price.

\begin{keyblock}
\textbf{Cell~6.2.13.} \emph{(Implied Volatility -- the rate-to-price
conversion.)}\\
Given the realized vol of the bond's YTM, $\sigma_y$, the corresponding
vol of the bond's price is approximately
\[
  \sigma_P \;\approx\; \mathrm{MD}(P) \cdot \sigma_y,
\]
where $\mathrm{MD}(P)$ is the modified duration of the bond at the
quote date. The vol of the \emph{forward} bond price is approximately
the same as $\sigma_P$ for our purposes (small adjustment for
forward-versus-spot price level).
\end{keyblock}

\begin{tipblock}
\textbf{Filling the gap.} The first-order linkage
$\Delta P / P \approx -\mathrm{MD} \cdot \Delta y$ is the modified-duration
identity. Squaring and taking expectations,
$\Var(\Delta P / P) \approx \mathrm{MD}^2 \, \Var(\Delta y)$, so
$\sigma_P \approx \mathrm{MD} \cdot \sigma_y$. This conversion ignores
convexity (a second-order term), which the lecturer flagged as
acceptable for a rough plug-in vol but not for a serious quote.
For the Freddie~Mac with $\mathrm{MD} \approx 7$--$10$ and
$\sigma_y \approx 2.68\%$ per year, $\sigma_P$ is somewhere in the
range of 19--27\% per year --- order-of-magnitude consistent with
liquid bond-option implied vols.
\end{tipblock}

\begin{checkpoint}
\textbf{My note.} The lecturer was explicit that this is a hack:
``we're really cheating, and we're saying we're just going to plug in
a realized vol for the implied vol.'' The case study notebook walks
the reader through the dollar pricing this produces. Lectures~6 and~7
will get into actual bond-option implied-vol surfaces (forward vol,
the vol skew); the same machinery will produce a more defensible
$\sigma$ input.
\end{checkpoint}

\section{Pricing the Freddie Mac callable end-to-end}\label{L5-sec:pricing}

With Black's general formula and a vol plug-in, the rest of the
pricing procedure is the bond-pricing-plus-option-subtraction logic
flagged earlier.

\subsection{The decomposition, operationally}\label{L5-sec:pricing-decomp}

\begin{keyblock}
\textbf{Cells~6.2.14--16.} \emph{(Pricing approach.)}\\
\[
  V^{\text{callable}}_t
   \;=\; V^{\text{vanilla}}_t \;-\; C_t.
\]
\begin{enumerate}[label=(\arabic*)]
  \item Build $V^{\text{vanilla}}_t$ from the discount curve $Z(t,T_i)$
        and the bond's coupon schedule:
        $V^{\text{vanilla}}_t = \sum_i c_i Z(t, T_i) + 100\, Z(t,
        T_{\text{bond}})$.
  \item Compute the forward price $F_t$ of the underlying bond \emph{to
        the option expiration} using the carry identity from
        Lecture~4.
  \item Plug $F_t$, $K=100$, $\sigma_P = \mathrm{MD}\cdot\sigma_y$, $Z(t,
        T_{\text{opt}})$, and $\tau = T_{\text{opt}} - t$ into Black's
        general call formula to get $C_t$.
  \item Subtract: $V^{\text{callable}}_t = V^{\text{vanilla}}_t - C_t$.
\end{enumerate}
\end{keyblock}

The lecturer emphasised that every input in this procedure has been
pre-built by earlier weeks of the course: $Z(\cdot, \cdot)$ comes
from Lecture~1's curve, $F_t$ from Lecture~4's forward-pricing
formula, the modified duration from Lecture~2's bond-math toolkit, and
the realized $\sigma_y$ from the time series of the YTM. ``This is why
I keep saying you need a good Python function to get forward prices
easily and to get that implied vol --- you'll use them every day for
the rest of the course.''

\subsection{Option-Adjusted Spread (OAS)}\label{L5-sec:oas}

The case study notebook shows the model price the procedure above
produces and compares it to the market quote. The two will not match
exactly. The OAS is the standardized way of expressing that gap.

\begin{keyblock}
\textbf{Cell~6.2.17.} \emph{(Option-Adjusted Spread.)}\\
The \emph{option-adjusted spread} is the parallel shift to the spot
discount curve that makes the model (here: Black's general formula)
reproduce the observed market price of the option-embedded security.
Formally, OAS is the unique scalar $\lambda$ such that
\[
  V^{\text{model}}_t\bigl(Z(t,\cdot)\,e^{-\lambda(\cdot - t)}\bigr)
  \;=\; V^{\text{market}}_t.
\]
A positive OAS means the security trades cheaper than the model would
predict given the curve --- it offers a yield pickup beyond what the
risk-free curve plus Black's formula explains.
\end{keyblock}

The lecturer made the parallel to YTM explicit again. ``YTM is just
reverse-engineered to match the market quote and help you understand
how expensive the bond is. OAS is just reverse-engineered to make
Black hit that number and tell you how expensive the option-embedded
security is.'' Two more practical points anchored why OAS is widely
quoted and not just a curiosity:

\begin{enumerate}[label=(\arabic*)]
  \item \emph{It is model-light.} OAS uses Black's general formula or a
        similarly standard model. Any market participant with the
        zero-coupon curve can compute it. There is no proprietary
        valuation engine required to read OAS off a Bloomberg screen,
        which is why ``Bloomberg defaults to showing OAS using
        Black, and sometimes Bachelier as well.''
  \item \emph{It compares apples to apples.} An OAS on a Freddie~Mac
        callable can be sized against an OAS on a callable-corporate or
        a mortgage-backed security. The dollar prices cannot be
        compared directly any more than dirty prices on different
        bonds can be compared.
\end{enumerate}

\begin{tipblock}
\textbf{Filling the gap.} The lecturer flagged what a non-zero OAS
\emph{does not} mean. ``If you truly believed Black's model is
appropriate not just for translation but for actual valuation, then any
OAS would be an arbitrage.'' But Black's formula is wrong in the same
ways for fixed-income options that it is wrong for equity options: the
implied-vol surface has a skew (Lecture~7), the underlying has fat
tails, and the replication argument fails in the presence of frictions
and finite liquidity. A 50~bp OAS on an illiquid Freddie~Mac is far
more likely a vol-skew effect or a liquidity premium than a free
lunch. ``A non-zero OAS is not a signal there's free money to be had,
but if it's substantially non-zero it would make us look twice.''
\end{tipblock}

\subsection{Negative convexity, visualised}\label{L5-sec:neg-convex}

The case study notebook produces a plot of model price as a function
of the level of rates: the vanilla bond on one axis (the standard
downward-sloping price-versus-rate curve, with positive convexity),
and the callable bond on the other (the downward-sloping bond curve
flattened from above by the call option's exercise barrier). The
lecturer used the picture to anchor the negative-convexity language.

\begin{keyblock}
\textbf{Cells~6.2.18--19.} \emph{(Differences with callable bonds.)}
\begin{itemize}
  \item The callable's price is \emph{lower} than the vanilla's at
        every rate level, by exactly the value of the embedded call.
  \item The callable price as a function of rates is \emph{negatively
        convex}: as rates fall, the call moves into the money for
        Freddie~Mac and the price stops rising near \$100. As rates
        rise, the bond loses value normally.
  \item The callable can rise modestly above \$100 between exercise
        dates --- the Bermudan structure leaves windows in which the
        issuer is mid-period and cannot call until the next semi-annual
        date.
\end{itemize}
\end{keyblock}

The economic interpretation is that the callable is a one-sided
instrument from the bondholder's perspective: capped on the upside,
not capped on the downside. ``You can lose a lot when rates go up, and
you can't gain when rates go down. What do we call that? Negatively
convex.''

\begin{checkpoint}
\textbf{My note.} The duration-hedging consequence is the practically
important one. A vanilla-bond holder who duration-hedges with swaps
benefits regardless of which way rates move (positive convexity earns
on volatility). A callable-bond holder who duration-hedges with swaps
\emph{loses} on rate volatility on either side --- the negative gamma
on the embedded call dominates the bond's positive convexity in
typical parameter ranges. The lecturer warned: ``if you bought this
callable bond and you duration-hedged it, you would essentially lose
either way if rates move, which is the opposite of our usual
result.''
\end{checkpoint}

\subsection{Are callable bonds systematically overpriced?}\label{L5-sec:longstaff}

\begin{keyblock}
\textbf{Cell~6.2.23.} \emph{(Are callable bonds overpriced?)}\\
Empirical work --- Longstaff and others --- has documented that
callable bonds in the data trade \emph{above} their model values
based on Black's formula. The standard narrative attributes this to
suboptimal exercise by issuers: the call option is held by an entity
(an agency, a corporate Treasury) for whom the dollar gain from
calling is small relative to the operational friction. So in the
data, calls happen later than they ``should'' happen, and the
embedded short call is therefore worth less to the bondholder than a
fully-optimal-exercise model would say.
\end{keyblock}

The lecturer pointed at Francis Longstaff (UCLA, of the Longstaff--Schwartz
American-option algorithm) as the canonical reference. The story the
data tells is that a model-implied price of \$95 against a market
price of \$98 has at least three potential explanations:
\begin{enumerate}[label=(\arabic*)]
  \item \emph{Wrong inputs.} The discount curve, the implied vol, the
        forward price are all noisy; small errors in any of them push
        the model price.
  \item \emph{Wrong model.} Black's formula assumes lognormal forward
        prices and instantaneous re-balancing. Bond options exhibit
        skew (Lecture~7) and the bonds themselves are illiquid enough
        that the replication argument is suspect.
  \item \emph{Suboptimal exercise.} The agency or corporate issuer
        does not call at the optimal time; the embedded call is worth
        less than a fully-rational-issuer model would say; subtracting
        a smaller call from the vanilla gives a higher callable price,
        consistent with the observed quote.
\end{enumerate}
The case study and exercise will revisit a real Freddie~Mac bond in
which exercise was visibly delayed. ``Even in the last two years of
doing this case study, I've seen that in the data --- they don't call
when it seems they should.''

\begin{exerciseblock}
\textbf{Exercise (assigned).} The case study notebook walks through
the full Black-based valuation of one of the Freddie~Mac callable
bonds in Cell~6.2.7's table, including: bootstrapping the discount
curve, computing $F_t$ to first call, plugging into Black's general
formula, and inferring the OAS implied by the market quote. A separate
exercise repeats the procedure for an option on a SOFR future, where
the underlying is a STIR future on an interest rate (not the rate
itself); compare to a Treasury-future option as a sanity check.
\end{exerciseblock}

\section{Looking ahead: caps, floors, and the time-mismatch quirk}\label{L5-sec:caps}

The lecturer had time to launch into the cap-and-floor primer that
properly belongs to Lecture~6's first section. He covered the
\emph{single-period} caplet and floorlet here; the multi-period
machinery (cap volatility, the term-structure of caplet vols) waits
for next week.

\subsection{Caplets and floorlets}\label{L5-sec:caplets}

\begin{keyblock}
\textbf{Slide L5.7.} \emph{(Caplet definition.)}\\
A caplet is a single-period European call option on an interest-rate
\emph{index} (e.g.\ SOFR, a 90-day T-bill rate). With strike $K$ (a
rate) and expiration $T$, the payoff at $T$ is
\[
  \text{payoff}_T \;=\; \frac{N}{m}\,\bigl(r_{\text{lock}} - K\bigr)^+,
\]
where $N$ is a notional, $m$ is the period frequency (4 for quarterly,
2 for semi-annual, 1 for annual), and $r_{\text{lock}}$ is the rate
\emph{set in arrears}: locked at $T - 1/m$ but paid at $T$.

A floorlet is the corresponding put: payoff
$\frac{N}{m}(K - r_{\text{lock}})^+$.
\end{keyblock}

The class flagged the timing convention as the most surprising feature.
A caplet on a quarterly rate that pays at $T = 1$~year locks in its
payoff at $T - 0.25 = 0.75$~years. ``All the uncertainty is resolved at
9~months, but you don't get paid for a year.'' The lecturer was
explicit that this is a \emph{convention} inherited from the LIBOR-era
floating leg of an interest-rate swap, designed to make the caplet a
clean hedge against a single floating-rate payment. ``It's mimicking
the way swaps' floating side has been quarterly for decades.''

\begin{tipblock}
\textbf{Filling the gap.} The lecturer pressed the unit-conversion
point: a strike of $K = 4\%$ and a locked rate of $r_{\text{lock}} =
4.5\%$ on a quarterly caplet pays $\frac{100}{4}(0.045 - 0.040) =
\$0.125$ per dollar of notional. The factor $1/m$ converts the
annualised rate-difference into a period-payment. The notional $N=
100$ is the same convention used throughout the bond chapter ---
keeping all dollar quantities scaled to face $100$ avoids unit
arithmetic mistakes when the same instrument is mixed with bonds and
swaps.
\end{tipblock}

\subsection{Caps and floors as portfolios; the time-mismatch}\label{L5-sec:caps-portfolio}

\begin{keyblock}
\textbf{Slide L5.8.} \emph{(Caps and floors.)}
\begin{itemize}
  \item A \emph{cap} is a portfolio of caplets on the same underlying
        rate, at the same strike $K$, with sequential expirations from
        a starting date out to the cap maturity $T$. ``A 10-year
        quarterly cap is 40 caplets all struck at $K$, with
        expirations one quarter, two quarters, $\ldots$, 40~quarters
        out.''
  \item The \emph{first} caplet (the one expiring at $T_1 = 1/m$) is
        \emph{excluded}: by the lock-in-arrears convention, its payoff
        is already known at $t=0$ (the rate has been observed), so
        it is a deterministic cash flow rather than an option.
  \item Black's general formula prices each caplet independently:
        \[
          c^{(j)}_t \;=\; \frac{N}{m}\,Z(t, T_j)
            \,\bigl[\, f^{(j)}_t\,\cN(d_1^{(j)})
                       \;-\; K\,\cN(d_2^{(j)})\,\bigr],
        \]
        with the forward rate $f^{(j)}_t$ for the period
        $[T_j - 1/m,\, T_j]$ and time-to-uncertainty-resolution
        $\tau^{(j)} = T_j - 1/m - t$ in $d_1^{(j)}, d_2^{(j)}$. The
        cap value is $\sum_j c^{(j)}_t$.
\end{itemize}
\end{keyblock}

The lecturer made a particular fuss about a subtlety that does not
arise for any other instrument in the course: the time appearing in
the discount factor $Z(t, T_j)$ is \emph{not} the same as the time
appearing in $\tau^{(j)}$ inside $d_1^{(j)}, d_2^{(j)}$. The discount
factor uses time to \emph{payment} ($T_j$) because that is when the
cash arrives. The variance term uses time to \emph{lock-in} ($T_j -
1/m$) because that is when the uncertainty resolves. ``If I were
giving a closed-book exam I would absolutely catch a bunch of wrong
answers for this.''

\begin{tipblock}
\textbf{Filling the gap.} Why is the variance term scaled by
$\tau^{(j)} = T_j - 1/m - t$? The implied vol $\sigma$ is the standard
deviation of the log forward rate over the period during which the
rate \emph{evolves}. Once the rate has been observed at $T_j - 1/m$,
its log returns have stopped accumulating, so the cumulative variance
is $\sigma^2 \tau^{(j)}$, not $\sigma^2(T_j - t)$. Equivalently:
between $T_j - 1/m$ and $T_j$ the payoff is a deterministic cash
amount, so its log-return variance is zero. Pick up the discount
factor for the gap $1/m$, drop the variance accrual for the gap.
\end{tipblock}

\begin{checkpoint}
\textbf{My note.} The mismatch between time-to-lock and time-to-pay is
a real arithmetic trap. Practical advice from the lecturer: write the
caplet pricing function with two separate time inputs, $T_{\text{pay}}$
and $T_{\text{lock}} = T_{\text{pay}} - 1/m$, and never reuse one for
the other. The discount factor $Z(t, T_{\text{pay}})$ scales the cash
flow; $\tau = T_{\text{lock}} - t$ scales the variance.
\end{checkpoint}

\subsection{Why caps from the swap curve?}\label{L5-sec:cap-curve}

A student asked the natural follow-up: should the discount factor
$Z(t, T_j)$ and the forward rate $f^{(j)}_t$ be bootstrapped from the
Treasury curve or from the swap curve? The lecturer gave a strong
recommendation.

\begin{keyblock}
\textbf{Slide L5.9.} \emph{(Curve choice for cap/floor pricing.)} For
caps, floors, and swaptions, the convention is to bootstrap discount
factors and forward rates from the \emph{swap curve}, not the
Treasury curve. The reasons:
\begin{itemize}
  \item Caps and floors are most often used to hedge swap floating
        legs, and pricing them on the swap curve makes the hedge
        relationship clean.
  \item The swap-spread story (\S\ref{L5-sec:swap-spread}) implies the
        swap curve and the Treasury curve disagree by tens of bps.
        Choosing the swap curve gives the right answer for swap-related
        derivatives.
  \item Discount factors are \emph{product-dependent} --- the same
        statement as the Lecture~4 case where the swap curve was
        bootstrapped separately from the Treasury curve.
\end{itemize}
\end{keyblock}

The lecturer's closing remark on this thread captured the recurring
theme of the course. ``I told you back in Week~1: you will not do a
single thing in this class where step one isn't get a curve that you
trust.'' Caps, floors, and swaptions are no exception, and the curve
that they trust is the swap curve.

\begin{checkpoint}
\textbf{My note.} The cap-pricing recipe is now five mechanical
steps, all of which use earlier course objects: (1)~build the swap
discount curve $Z^{\text{swap}}(t, \cdot)$ using the bootstrap from
Lecture~4; (2)~extract forward rates $f^{(j)}_t$ from the swap curve;
(3)~price each caplet via Black's general formula with two distinct
time inputs ($T_j$ for the discount, $T_j - 1/m$ for the variance);
(4)~scale by $N/m$; (5)~sum across $j$, excluding the first caplet.
The remaining hard input --- the implied vol surface for caplets ---
is what Lecture~6 introduces under the heading ``forward
volatility.''
\end{checkpoint}

The lecturer closed with a quick announcement that homework~1 will
exercise this exact procedure on a small cap and a small floor and
``build us up to the actual hard problem next week.''

\section{Closing the loop}\label{L5-sec:close}

The lecture covered what is, structurally, a small set of changes to
FINMAT~330's option-pricing toolkit:

\begin{enumerate}[label=(\arabic*)]
  \item Replace the spot-price form of Black--Scholes with the
        forward-price form. This absorbs carry and makes the
        underlying nearly Brownian even for instruments (bonds,
        rates) whose spot dynamics are obviously not.
  \item Replace the constant-rate discount $e^{-r(T-t)}$ with the
        observable discount factor $Z(t,T)$. The technical justification
        --- the change of numeraire to the $T$-forward measure ---
        lives in FINMAT~330; the operational consequence is that the
        same curve-bootstrapping you already do for cash bonds
        produces the inputs for option pricing.
  \item Treat implied volatility as a quoting language, not as a
        belief. Black-implied vol is to options what YTM is to
        bonds: a single comparable scale across strikes, expirations,
        and exercise styles. Even American and Bermudan options are
        quoted in Black-implied vol.
  \item For our specific Freddie~Mac problem, value the callable as
        $V^{\text{vanilla}} - C$, with the call $C$ priced via Black's
        general formula and a vol input approximated by
        $\mathrm{MD}\cdot\sigma_y$. Compare to the market quote via
        OAS, but do not over-interpret a non-zero OAS.
  \item For caps and floors --- which are the marquee instruments for
        Lecture~6 onward --- price each caplet/floorlet with Black's
        general formula and the lock-in-arrears time mismatch, using
        the swap curve for both the discount factors and the forward
        rates.
\end{enumerate}

The Freddie~Mac case study and an exercise on options on STIR futures
are the working homework for this material. The next lecture begins
the pivot to forward volatility, which underpins serious cap, floor,
and swaption pricing --- replacing the realized-vol-times-modified-duration
hack used here with a proper implied-vol surface.
